\section{Experimental design and statistical analysis}
\label{sec:design}
\subsection{Sequential controls and preserved stages}
The historical ladder changes one broad design component at a time, while
explicit controls expose remaining differences. B0-UW addresses weighting;
50-round no-noise controls address undertraining; E1 clipping-only controls
separate clipping from additional perturbation; and the private popularity
reference tests whether iterative personalization is competitive at all.
No single comparison isolates every causal mechanism. In particular,
parameterization changes optimization and sensitivity geometry together.

\begin{table}[tbp]\centering\small
\caption{Experimental stages and their evidence units.}
\label{tab:stages}
\begin{tabularx}{\linewidth}{lYY}\toprule
Stage & Question & Retained evidence \\\midrule
Phase 0 & Can the ranking protocol be reproduced? & Data fingerprints, evaluator checks, random and raw-popularity references. \\
B0--B4 & How do weighting, federation, privacy, and rank interact? & Three-seed weighting control; five-seed final federated comparisons. \\
E1 & Does a fixed public factor help under equal tuning? & 648 search jobs; 270 final collaborative states; 21 popularity arrays. \\
E2/E2b & Is local-state disturbance persistent and coupling-sensitive? & 120 pulse/reset probes and 40 aligned follow-up probes from 30 checkpoints. \\
E3 & Does the diagnostic replicate and depend on balancing? & 30 ML-1M training runs; 200 controlled probes across two datasets. \\
E4 & Do non-Gaussian margins matter in trained states? & 50 states; 3,200 prescribed item pairs; six synthetic cases. \\
\bottomrule\end{tabularx}
\tabnote{A probe is not an independently trained model. Some E1 final states reuse
eligible search runs, and E2--E4 deliberately reuse frozen checkpoints.}
\end{table}

E1's grid has $9\times4\times3\times6=648$ search jobs, including eleven
retained failures. The selected units, six levels, and five seeds yield
$9\times6\times5=270$ finite final states. Reuse of eligible search cells
gives 756 unique training records rather than 918. Twenty noisy popularity
arrays and one deterministic bounded no-noise array bring one-time extension
test scoring to 291 models. Selections and validation-chosen ranks were frozen
at 15:37:41 UTC on 2 October 2026, before the first extension test scoring.
This freeze is a stage-specific safeguard, not evidence that the dataset had
never been examined earlier.

\subsection{Seeds and paired uncertainty}
The centralized weighting comparison uses seeds 42, 123, and 2026. Final
B1--B4 and E1 use those plus 7 and 99. Utility tables give the arithmetic mean
and sample SD across the retained seeds. To compare two methods on NDCG@10,
first average the metric over seeds separately for each user, then form the
paired user difference. A bootstrap samples users with replacement from
these paired differences 100,000 times with bootstrap seed 2026.

The primary B3 family contains four comparisons against B1; B4 contains
sixteen rank-by-privacy comparisons against B2; E1 contains sixteen matched-
rank FixedB-minus-Two comparisons. Bonferroni-adjusted interval levels are
98.75\% for B3 and 99.6875\% for each sixteen-comparison family. These are
family-wise interval constructions conditional on the saved training runs
and population represented by the benchmark. They do not jointly quantify
uncertainty over new datasets and independently sampled training procedures.

An adjusted interval entirely above zero supports a positive comparison within
its family. An interval containing zero is recorded as ``includes zero,''
not as proof of equivalence. Ordinary 95\% intervals for secondary comparisons
are retained in the supplement but do not inherit the primary multiplicity
interpretation. Figure error bars are seed SDs, and their overlap is not the
criterion used for primary decisions.

\subsection{Rank selection, groups, and personalization diagnostics}
For presentation of a single low-rank curve, each method's rank is selected
from validation at the relevant noise level. That curve is different from
the matched-rank primary comparison. A test-best rank never replaces a
validation-selected rank because it happens to look better in the final table.
The historical B4 and E1 selections differ and are retained separately.

Activity and cold-target tables are descriptive. Their groups are not
independent discoveries, and the fourteen cold targets are too few for a
strong cold-item claim. E1 also compares validation predictions using each
user's own vector, permuted vectors, and a mean vector. These 810 retained
rows diagnose the contribution of local personalization. They are oracle
analyses of stored states, not alternative serving procedures with analyzed
privacy guarantees.

\subsection{Diagnostic replication and advancement criteria}
E2 averages two pulse replicates within a checkpoint seed before summarizing
across five seeds. E2b follows the same rule. Their means, medians, ranges,
and SDs are descriptive; correlated users, intervention branches, and draws
are not treated as independent training runs. E3 keeps transferred training
settings fixed and uses a preselected observer population. E4 declares its
checkpoint set, pair types, and practical-error gate before examining outputs.
The gate requires absolute matched-Gaussian probability error greater than
$0.01$ on at least 5\% of primary pairs in both datasets. It is a research
screening criterion, not a significance threshold.

Protocols, source hashes, checkpoint identities, failure inventories, and
observations are retained. This separation of initial evidence, adaptive
follow-ups, and final decisions makes the chronological research process
auditable without pretending every later question was specified at the outset.
